% All the definitions and results that Section~\ref{sec:fs-hints} and
% Section~\ref{sec:eufcma-gen} build on are collected here, so that
% extended_proof.tex can be read without the companion paper at hand.  Every
% imported statement carries, in its name, the number it has in the source it
% is taken from, which is the number that extended_proof.tex cites.

\subsection{Notation}\label{sec:notation}

We let $\secp \in \NN$ be the security parameter; all algorithms implicitly take
the unary encoding $1^\secp$ as an input.
We write PPT for ``probabilistic polynomial time in the length of the inputs'',
$\poly(x_1 \until x_n)$ for an unspecified positive polynomial and
$\negl(\secp)$ for an unspecified negligible function of $\secp$.
When an algorithm $\adv$ has black-box query access to an oracle $\O$ we write
$\adv^{\O}$.
For $x, y \in \set{0,1}^*$ we write $x \concat y$ for their concatenation.

For a probability distribution $D$ we write $x \gets D$ for sampling $x$
from $D$; for a finite set $S$ we write $\unifover{S}$ for the uniform
distribution over $S$ and $x \sample S$ for sampling from $\unifover{S}$.
The statistical distance between two distributions $D_0, D_1$ over a finite set
$S$ is
\begin{equation}
    \statdist{D_0}{D_1} \defeq
    \frac{1}{2} \sum_{s \in S} \abs{\Pr[D_0 = s] - \Pr[D_1 = s]},
\end{equation}
and for an algorithm $\adv$ outputting a bit we write
\begin{equation}
    \distadv{D_0}{D_1}{\adv} \defeq
    \abs{\Pr[1 \gets \adv(x) \mid x \gets D_0] - \Pr[1 \gets \adv(x) \mid x \gets D_1]}
\end{equation}
for its advantage in distinguishing $D_0$ from $D_1$.
Recall that $\distadv{D_0}{D_1}{\adv} \leq \statdist{D_0}{D_1}$ for every
$\adv$, with equality for an unbounded $\adv$~\cite[Proposition~2.3]{C:ABDPW25}.

\subsection{Relations and \texorpdfstring{$\Sigma$}{Sigma}-protocols}\label{sec:sigma}

A binary relation $R$ is a set of tuples $(x,w) \in \cX \times \cW \subseteq
\set{0,1}^* \times \set{0,1}^*$; we call $x$ the \emph{statement} and $w$ the
\emph{witness}.
All relations we consider are NP relations and come equipped with an
\emph{instance generator} $\Gen_R$, a PPT algorithm which on input $1^\secp$
outputs a pair $(x,w) \in R$ with $\abs{x} = \poly(\secp)$.

\begin{definition}[Hard relation, {\cite[Definition~2.1]{C:ABDPW25}}]\label{def:hard-rel}
    A relation $R$ is a \emph{hard relation} if for all PPT algorithms $\adv$
    \begin{equation}
        \reladv{R}{\adv} \defeq
        \Pr\left[(x, w^*) \in R \;\middle|\;
            (x,w) \gets \Gen_R(1^\secp),\ w^* \gets \adv(x)\right]
        = \negl(\secp).
    \end{equation}
\end{definition}

\begin{definition}[$\Sigma$-protocol, {\cite[Definition~2.2]{C:ABDPW25}}]\label{def:sigma}
    A \emph{$\Sigma$-protocol} for a relation $R$ is a $3$-message interactive
    protocol between a prover and a verifier, given by a tuple of PPT algorithms
    $\Sigma = (\Prover_1, \Prover_2, \Verif)$, as follows.
    \begin{enumerate}
        \item On input $(x,w) \in R$ the prover runs $(\com, \state) \gets
            \Prover_1(x,w)$, obtaining a \emph{commitment} $\com$ and the state
            $\state$ needed for the response, and sends $\com$ to the verifier.
        \item The verifier samples a \emph{challenge} $\chall \sample \cC$
            uniformly from the challenge space $\cC$ and sends it to the prover.
        \item The prover replies with the \emph{response} $\rsp \gets
            \Prover_2(\state, \chall)$.
        \item The verifier runs $b \gets \Ver_\Sigma(x, \com, \chall, \rsp)$ and
            accepts if and only if $b = 1$.
    \end{enumerate}
    We always assume \emph{correctness}: an honest execution is accepted except
    with probability $\negl(\secp)$.
\end{definition}

The protocol is \emph{weak honest-verifier zero-knowledge} (\wHVZK) if there is
a PPT simulator producing transcripts $(\com, \chall, \rsp)$ that no PPT
algorithm can distinguish from honest ones, the challenge being uniform in both
cases~\cite[Definition~2.3]{C:ABDPW25}; \Cref{def:hint-hvzk} below is the
hint-assisted relaxation of this notion that we shall actually use.

\begin{definition}[MinEnt, {\cite[Definition~2.4]{C:ABDPW25}}]\label{def:minent}
    $\Sigma$ has \emph{high commitment min-entropy} if
    \begin{equation}
        \minentropy{\Sigma} \defeq \max_{(x,w)} \max_{\com}
        \Pr\left[\com = \com' \;\middle|\; (\com', \state) \gets \Prover_1(x,w)\right]
        = \negl(\secp),
    \end{equation}
    where the first maximum ranges over the pairs that $\Gen_R(1^\secp)$ may
    output.
\end{definition}

Note that $\minentropy{\Sigma}$ is a probability, not a number of bits: the
actual min-entropy of the commitment is $-\log(\minentropy{\Sigma})$.

\begin{definition}[Special soundness, {\cite[Definition~2.5]{C:ABDPW25}}]\label{def:ss}
    $\Sigma$ is \emph{special sound} if, given a statement $x$ and two accepting
    transcripts $(\com, \chall, \rsp)$ and $(\com, \chall', \rsp')$ for $x$ with
    $\chall \neq \chall'$, one can compute in time $\poly(\abs{x})$ a witness
    $w$ such that $(x,w) \in R$.
\end{definition}

For \sqisign{}, and more generally for the protocols we are interested in,
special soundness only holds with respect to a \emph{different} relation, which
is the reason for the following relaxation.

\begin{definition}[Soundness relation, {\cite[Definition~3.4]{C:ABDPW25}}]\label{def:sound-rel}
    Let $\Sigma$ be a $\Sigma$-protocol for the relation $R$.
    A relation $\tilde{R} \subseteq \tilde{\cX} \times \tilde{\cW}$ is
    \emph{compatible} with $R$ if $\cX = \tilde{\cX}$ and their instance
    generators sample statements with the same distribution.
    We say that $\Sigma$ is \emph{special sound with respect to a compatible
    relation $\tilde{R}$} if from a statement $x$ and two accepting transcripts
    $(\com, \chall, \rsp)$, $(\com, \chall', \rsp')$ with $\chall \neq \chall'$
    one can compute in time $\poly(\abs{x})$ a witness $w^*$ with $(x, w^*) \in
    \tilde{R}$.
    In this case $\tilde{R}$ is called a \emph{soundness relation} for $\Sigma$.
\end{definition}

A $\Sigma$-protocol can be read as an identification scheme, the prover trying
to convince the verifier that it knows a witness for the identity $x$.
The relevant security notion is security against \emph{passive impersonation
attacks}: the attacker may first eavesdrop on honest executions, modelled by the
transcript oracle $\OTrans$ of \Cref{fig:imppa}, and then tries to impersonate
the prover.

\begin{figure}[htb]
    \centering
    \begin{pchstack}
        \procedure[linenumbering]{$\imppatxt_\Sigma(\adv)$}{
            (x, w) \gets \Gen_R(1^\secp) \\
            (\com^*, \state^*) \gets \adv_1^{\OTrans_{x,w}}(x) \\
            \chall^* \sample \cC \\
            \rsp^* \gets \adv_2(\state^*, \chall^*) \\
            \pcreturn \Ver_\Sigma(x, \com^*, \chall^*, \rsp^*)
        }
        \pchspace
        \procedure[linenumbering]{$\OTrans_{x,w}()$}{
            (\com, \state) \gets \Prover_1(x, w) \\
            \chall \sample \cC \\
            \rsp \gets \Prover_2(\state, \chall) \\
            \pcreturn (\com, \chall, \rsp)
        }
    \end{pchstack}
    \caption{The $\imppatxt_\Sigma$ game of~\cite[Game~1]{C:ABDPW25}, played by a
        two-stage adversary $\adv = (\adv_1, \adv_2)$.}\label{fig:imppa}
\end{figure}

\begin{definition}[\IMPPA, {\cite[Definition~2.7]{C:ABDPW25}}]\label{def:imppa}
    $\Sigma$ is secure under passive impersonation attacks (\IMPPA) if for all
    PPT two-stage algorithms $\adv = (\adv_1, \adv_2)$
    \begin{equation}
        \imppaadv{\Sigma}{\adv} \defeq \Pr[\imppatxt_\Sigma(\adv) = 1]
        = \negl(\secp),
    \end{equation}
    where $\imppatxt_\Sigma$ is the game of \Cref{fig:imppa}.
\end{definition}

\subsection{Signature schemes in the ROM}\label{sec:sig-rom}

In the random oracle model (ROM) all algorithms are given black-box query access
to the random oracle $\RO$, implementing a uniformly random function
$\set{0,1}^* \to \cC$; we usually omit the dependency on $\RO$ from the
notation.

\begin{definition}[Signature scheme, {\cite[Definition~2.8]{C:ABDPW25}}]\label{def:sigscheme}
    A \emph{signature scheme} in the ROM for a message space $M$ is a tuple of
    PPT algorithms $\Sigscheme = (\Gen, \Sign, \Ver)$ where $\Gen(1^\secp)$
    outputs a key pair $(\pk, \sk)$, $\Sign(\sk, \msg)$ outputs a signature
    $\sig$ for $\msg \in M$, and $\Ver(\pk, \msg, \sig)$ outputs $1$ (accept) or
    $0$ (reject).
    Correctness requires that honestly generated signatures verify except with
    probability $\negl(\secp)$.
\end{definition}

\begin{figure}[htb]
    \centering
    \begin{pchstack}
        \procedure[linenumbering]{$\EUFCMA_\Sigscheme(\adv)$}{
            (\pk, \sk) \gets \Gen(1^\secp) \\
            \Qset \coloneqq \emptyset \\
            (\msg^*, \sig^*) \gets \adv^{\OSign}(\pk) \\
            \pcif \Ver(\pk, \msg^*, \sig^*) = 1 \wedge \msg^* \notin \Qset \pcthen \pcreturn 1 \\
            \pcreturn 0
        }
        \pchspace
        \procedure[linenumbering]{$\OSign(\msg)$}{
            \sig \gets \Sign(\sk, \msg) \\
            \Qset \coloneqq \Qset \cup \set{\msg} \\
            \pcreturn \sig
        }
    \end{pchstack}
    \caption{The $\EUFCMA_\Sigscheme$ game of~\cite[Game~2]{C:ABDPW25}.}\label{fig:eufcma}
\end{figure}

\begin{definition}[\EUFCMA, {\cite[Definition~2.9]{C:ABDPW25}}]\label{def:eufcma}
    A signature scheme $\Sigscheme$ satisfies \emph{existential unforgeability
    under chosen message attacks} (\EUFCMA) if for all PPT random-oracle
    algorithms $\adv$
    \begin{equation}
        \eufcmaadv{\Sigscheme}{\adv} \defeq
        \Pr[\EUFCMA_\Sigscheme(\adv) = 1] = \negl(\secp),
    \end{equation}
    where $\EUFCMA_\Sigscheme$ is the game of \Cref{fig:eufcma}.
\end{definition}

\begin{definition}[Fiat--Shamir signature, {\cite[Definition~2.10]{C:ABDPW25}}]\label{def:fs}
    Let $\Sigma$ be a $\Sigma$-protocol for a relation $R$ and let $M \subseteq
    \set{0,1}^*$.
    The signature scheme $\SIG[\Sigma]$ is defined as follows.
    \begin{itemize}
        \item $\Gen(1^\secp)$: sample $(x,w) \gets \Gen_R(1^\secp)$, set $\pk
            \coloneqq x$ and $\sk \coloneqq (x,w)$.
        \item $\Sign(\sk, \msg)$: generate a transcript $(\com, \chall, \rsp)$
            with the prover algorithms of $\Sigma$, but computing the challenge
            as $\chall \gets \RO(\com \concat x \concat \msg)$; output $\sig
            \coloneqq (\com, \chall, \rsp)$.
        \item $\Ver(\pk, \msg, \sig)$: accept if and only if $\sig = (\com,
            \chall, \rsp)$ with $\Ver_\Sigma(x, \com, \chall, \rsp) = 1$ and
            $\chall = \RO(\com \concat x \concat \msg)$.
    \end{itemize}
\end{definition}

That is, $\SIG[\Sigma]$ is the Fiat--Shamir transform $\FS{\Sigma}$ of $\Sigma$,
in which the challenge is obtained as $\chall \gets \RO(\com \concat x \concat
\msg)$, with the relation modified to
\begin{equation}
    R^* = \set{(x \concat \msg, w) \;\middle|\; (x,w) \in R,\ \msg \in M}.
\end{equation}
The standard reduction of its unforgeability to the \IMPPA security of $\Sigma$
is the result whose proof Section~\ref{sec:fs-hints} reorganises.

\begin{lemma}[Fiat--Shamir from \IMPPA, {\cite[Lemma~3.5]{EC:AABN02}}, {\cite[Lemma~2.2]{C:ABDPW25}}]\label{lem:aabn}
    For any PPT algorithm $\adv$ attacking the \EUFCMA security of
    $\SIG[\Sigma]$ there exists a PPT two-stage algorithm $\bdv$ such that
    \begin{equation}\label{eq:aabn}
        \eufcmaadv{\SIG[\Sigma]}{\adv} \leq
        (\numRO + 1) \cdot \imppaadv{\Sigma}{\bdv}
        + (\numRO + \numOSign + 1)\,\numOSign \cdot \minentropy{\Sigma},
    \end{equation}
    where $\numRO$ and $\numOSign$ are upper bounds on the number of queries
    that $\adv$ makes to $\RO$ and $\OSign$ respectively.
    Furthermore, $\bdv$ makes at most $\numOSign$ queries to the transcript
    oracle $\OTrans$.
\end{lemma}

\subsection{Fiat--Shamir with hints}\label{sec:fs-with-hints}

The framework of~\cite{C:ABDPW25} captures the protocols, such as \sqisign{},
whose simulator needs some additional information --- a \emph{hint} --- in order
to produce a transcript.

\begin{definition}[Hint distribution, {\cite[Definition~3.1]{C:ABDPW25}}]\label{def:hint-distr}
    Let $R \subseteq \cX \times \cW$ be a relation.
    A \emph{hint distribution} $\hintdistr$ for $R$ is a collection of
    distributions $\hintdistr = \set{\hintdistr_x}_{x \in \cX}$, where
    $\hintdistr_x \colon \HintSet_x \rightarrow [0,1]$ and the hints of $\HintSet_x$
    are efficiently representable in $\abs{x}$.
    The distribution $\hintdistr_x$ need not be efficiently sampleable.
    We call $(R, \hintdistr)$ a \emph{relation with hints}.
\end{definition}

\begin{figure}[htb]
    \centering
    \begin{pchstack}
        \procedure[linenumbering]{$\Real_\Sigma(1^\secp, \hintnum)$}{
            (x, w) \gets \Gen_R(1^\secp) \\
            \pcfor i = 1 \until \hintnum \pcdo \\
            \pcind (\com_i, \state_i) \gets \Prover_1(x, w) \\
            \pcind \chall_i \sample \cC \\
            \pcind \rsp_i \gets \Prover_2(\state_i, \chall_i) \\
            \pcind \pi_i \coloneqq (\com_i, \chall_i, \rsp_i) \\
            \pcreturn (x, \pi_1 \until \pi_\hintnum)
        }
        \pchspace
        \procedure[linenumbering]{$\HintSim_\Sigma(1^\secp, \hintnum, \Sim)$}{
            (x, w) \gets \Gen_R(1^\secp) \\
            \pcfor i = 1 \until \hintnum \pcdo \\
            \pcind \hint_i \gets \hintdistr_x \\
            \pcind \pi_i \gets \Sim(x, \hint_i) \\
            \pcreturn (x, \pi_1 \until \pi_\hintnum)
        }
    \end{pchstack}
    \caption{The hint-assisted \wHVZK experiment of~\cite[Experiment~2]{C:ABDPW25}.}
    \label{fig:hint-hvzk}
\end{figure}

\begin{definition}[Hint-assisted \wHVZK, {\cite[Definition~3.2]{C:ABDPW25}}]\label{def:hint-hvzk}
    Let $\Sigma$ be a $\Sigma$-protocol for a relation $R$ and let $\hintdistr$
    be a hint distribution for $R$.
    We say that $\Sigma$ is \emph{$\hintdistr$-hint-assisted \wHVZK} if there
    exists a PPT algorithm $\Sim$, the \emph{$\hintdistr$-hint-assisted
    simulator}, such that for all $\hintnum = \poly(\secp)$ and all PPT
    algorithms $\ddv$
    \begin{equation}
        \hinthvzkadv{\Sigma}{\hintdistr}{\Sim}{\ddv, \hintnum} \defeq
        \distadv{\Real_\Sigma(1^\secp, \hintnum)}%
                {\HintSim_\Sigma(1^\secp, \hintnum, \Sim)}{\ddv}
        = \negl(\secp),
    \end{equation}
    where $\Real_\Sigma$ and $\HintSim_\Sigma$ are the distributions defined in
    \Cref{fig:hint-hvzk}.
\end{definition}

Since the reduction hands the hints to the adversary, the relation has to remain
hard in their presence.

\begin{definition}[Hard relation with hints, {\cite[Definition~3.3]{C:ABDPW25}}]\label{def:hard-rel-hints}
    Let $R$ be a hard relation and let $\hintdistr$ be a hint distribution for
    $R$.
    The pair $(R, \hintdistr)$ is a \emph{hard relation with hints} if for all
    PPT algorithms $\adv$ and all $\hintnum = \poly(\secp)$
    \begin{equation}
        \hintreladv{R}{\hintdistr}{\adv}{\hintnum} \defeq
        \Pr\left[(x, w^*) \in R \;\middle|\;
            \begin{aligned}
                & (x,w) \gets \Gen_R(1^\secp), \\
                & \hint_1 \until \hint_\hintnum \gets \hintdistr_x, \\
                & w^* \gets \adv(x, \hint_1 \until \hint_\hintnum)
            \end{aligned}
        \right] = \negl(\secp).
    \end{equation}
\end{definition}

The intermediate notion through which the reduction passes is a
\emph{non-interactive} variant of \IMPPA: the adversary no longer has a
transcript oracle, but is given hints from which it can simulate transcripts
itself.

\begin{figure}[htb]
    \centering
    \procedure[linenumbering]{$\hintdistr\text{-}\hintimppatxt_\Sigma(\adv, \hintnum)$}{
        (x, w) \gets \Gen_R(1^\secp) \\
        \hint_1 \until \hint_\hintnum \gets \hintdistr_x \\
        (\com, \state) \gets \adv_1(x, \hint_1 \until \hint_\hintnum) \\
        \chall \sample \cC \\
        \rsp \gets \adv_2(\state, \chall) \\
        \pcreturn \Ver_\Sigma(x, \com, \chall, \rsp)
    }
    \caption{The $\hintdistr$-$\hintimppatxt_\Sigma$ game
        of~\cite[Game~3]{C:ABDPW25}.}\label{fig:hint-imppa}
\end{figure}

\begin{definition}[\hintIMPPA, {\cite[Definition~3.5]{C:ABDPW25}}]\label{def:hint-imppa}
    Let $(R, \hintdistr)$ be a relation with hints and let $\Sigma$ be a
    $\Sigma$-protocol for $R$.
    We say that $\Sigma$ is secure under $\hintdistr$-hint-assisted passive
    impersonation attacks ($\hintdistr$-\hintIMPPA) if for all $\hintnum =
    \poly(\secp)$ and all PPT two-stage algorithms $\adv = (\adv_1, \adv_2)$
    \begin{equation}\label{eq:hint-imppa-adv}
        \hintimppaadv{\Sigma}{\hintdistr}{\adv}{\hintnum} \defeq
        \Pr\left[\hintdistr\text{-}\hintimppatxt_\Sigma(\adv, \hintnum) = 1\right]
        = \negl(\secp),
    \end{equation}
    where $\hintdistr$-$\hintimppatxt_\Sigma$ is the game of
    \Cref{fig:hint-imppa}.
\end{definition}

Note that, unlike in the \IMPPA game, $\adv$ has \emph{no} access to a
transcript oracle, and that this game is not in the random oracle model: an
adversary that needs a random oracle simulates it internally.

The three results below are the ones that Section~\ref{sec:fs-hints} and
Section~\ref{sec:eufcma-gen} reorganise and reuse.
\Cref{lem:hint-imppa} is the step that trades the transcript oracle for the
hint-assisted simulator, and \Cref{lem:ss-hints} the step that concludes by
special soundness.

\begin{lemma}[Hint-assisted transcripts, {\cite[Lemma~3.1]{C:ABDPW25}}]\label{lem:hint-imppa}
    Let $(R, \hintdistr)$ be a relation with hints, let $\Sigma$ be a
    $\Sigma$-protocol for $R$ and let $\Sim$ be a $\hintdistr$-hint-assisted
    simulator for $\Sigma$.
    For any two-stage PPT algorithm $\adv = (\adv_1, \adv_2)$ against the
    \IMPPA security of $\Sigma$ there exist a two-stage PPT algorithm $\bdv$ and
    a PPT algorithm $\ddv$ such that
    \begin{equation}
        \imppaadv{\Sigma}{\adv} \leq
        \hintimppaadv{\Sigma}{\hintdistr}{\bdv}{\hintnum}
        + \hinthvzkadv{\Sigma}{\hintdistr}{\Sim}{\ddv, \hintnum},
    \end{equation}
    where $\hintnum$ is an upper bound on the number of queries that $\adv_1$
    makes to $\OTrans$.
\end{lemma}

\begin{lemma}[Special soundness with hints, {\cite[Lemma~3.2]{C:ABDPW25}}]\label{lem:ss-hints}
    Let $(R, \hintdistr)$ be a relation with hints and let $\Sigma$ be a
    $\Sigma$-protocol for $R$ with challenge space $\cC$ that is special sound
    with respect to a soundness relation $\tilde{R}$.
    For any two-stage PPT algorithm $\adv = (\adv_1, \adv_2)$ playing the
    $\hintIMPPA$ game for $\Sigma$ with $\hintnum$ hints from $\hintdistr$,
    there exists an expected polynomial time algorithm $\cE$ playing the game
    for $(\tilde{R}, \hintdistr)$ such that
    \begin{equation}
        \hintimppaadv{\Sigma}{\hintdistr}{\adv}{\hintnum} \leq
        \hintreladv{\tilde{R}}{\hintdistr}{\cE}{\hintnum} + \frac{1}{\abs{\cC}}.
    \end{equation}
    The expected running time of $\cE$ is about the same as that of the
    knowledge extractor of $\Sigma$, which invokes $\adv$ twice in expectation.
\end{lemma}

\begin{corollary}[Expected to strict polynomial time, {\cite[Corollary~3.1]{C:ABDPW25}}]\label{cor:exp-to-ppt}
    Let $\adv$ be an algorithm that runs in expected time $t$ with success
    probability $\eps$.
    Then there exists an algorithm $\bdv$ that runs in time at most $2t/\eps$
    and succeeds with probability at least $\eps/2$.
    In particular, if $t$ and $\eps^{-1}$ are polynomial in $\secp$, then $\bdv$
    runs in polynomial time in $\secp$ and has non-negligible success
    probability.
\end{corollary}

Chaining \Cref{lem:aabn}, \Cref{lem:hint-imppa} and \Cref{lem:ss-hints} gives the
main theorem of the framework, which Theorem~\ref{thm:EUFCMA-gen} improves upon.

\begin{theorem}[hint-rel reduces to \EUFCMA, {\cite[Theorem~1]{C:ABDPW25}}]\label{thm:abdpw}
    Let $(R, \hintdistr)$ be a relation with hints.
    Let $\Sigma$ be a $\Sigma$-protocol for $R$ that has challenge space $\cC$
    and is special sound with respect to a soundness relation $\tilde{R}$, and
    let $\Sim$ be a $\hintdistr$-hint-assisted simulator for $\Sigma$.
    For any PPT algorithm $\adv$ against the \EUFCMA security of $\SIG[\Sigma]$
    there exist a PPT algorithm $\bdv$ and an expected polynomial time algorithm
    $\cE$ such that
    \begin{equation}\label{eq:abdpw}
        \begin{aligned}
            \eufcmaadv{\SIG[\Sigma]}{\adv}
            &\leq (\numRO+1)\cdot
            \left(\hintreladv{\tilde{R}}{\hintdistr}{\cE}{\numOSign}
            + \hinthvzkadv{\Sigma}{\hintdistr}{\Sim}{\bdv, \numOSign}
            + \frac{1}{\abs{\cC}}\right) \\
            &\phantom{\leq ~}
            + (\numRO + \numOSign + 1)\,\numOSign \cdot \minentropy{\Sigma},
        \end{aligned}
    \end{equation}
    where $\numRO$ and $\numOSign$ are upper bounds on the number of queries
    that $\adv$ makes to $\RO$ and $\OSign$ respectively.
\end{theorem}

\subsection{Game-playing proofs}\label{sec:gameops}

For a game $\bG$ we write $\bG^\adv \Rightarrow 1$ for the event that the game
outputs $1$ when run with $\adv$, and for $4 \leq i \leq 6$ we let $\goodev_i$ be
the event that the $i$-th game never sets $\badev$.
Recall that two games are \emph{identical-until-$\badev$} if their code differs
only in statements that follow the setting of the flag to \btrue, and recall the
two standard consequences of this~\cite{EC:BelRog06}:
\begin{itemize}
    \item for identical-until-$\badev$ games $\bG, \bG'$ and any adversary
        $\adv$,
        \begin{equation}\label{eq:fundamental}
            \abs{\Pr[\bG^\adv \Rightarrow 1] - \Pr[\bG'^\adv \Rightarrow 1]}
            \leq \Pr[\bG^\adv \text{ sets the flag}],
        \end{equation}
        where the right-hand side may equally be taken to be $\Pr[\bG'^\adv
        \text{ sets the flag}]$, since identical-until-$\badev$ games set the
        flag with the same probability;
    \item and, denoting by $\goodev, \goodev'$ the events that the flag is never
        set,
        \begin{equation}\label{eq:good}
            \Pr[\bG^\adv \Rightarrow 1 \wedge \goodev]
            = \Pr[\bG'^\adv \Rightarrow 1 \wedge \goodev'].
        \end{equation}
\end{itemize}
